\scanpage{22}
\mspage{\handwritten{13}}
\uline{Remarques} 7.4. a) Remords à 7.3. Inutile de supposer $n$ pair et $n\handwritten{\geq}4$, en énonçant la conclusion sous la forme: le niveau de $E^{n-1}(Y_t)$ est égal à $n-1$ $(n\handwritten{\geq}1)$. On n'a pas besoin en effet d'un c\unsure{r|\deleted{t}r|\deleted{l}r}itère de Ogg-Néron-Chafarévitch ici, car il suffit de savoir que \uline{si} $E^{n-1}(Y_t)$ est de niveau $\handwritten{<}n-1$, i.e. s'il existe sur \deleted{$k(t)$} $\bar{k}'$ $(k'=k(t))$ un cycle sur un produit $Z\mathbin{\mathrm{x}}Y_{\overline{k'}}$ définissant un \deleted{iso}\inserted{épi}morphisme $H^{n-3}(Z,\underline{Q}_{\ell})(-1)\handwritten{\longrightarrow}H^{n-1}(Y_{\bar{k}'},\underline{Q}_{\ell})\handwritten{\longrightarrow}E(Y_{\bar{k}'}$ \deleted{alors, \unsure{quitte à passer} pour} alors pour tout $s\handwritten{\in}D$ assez voisin de $t$, la même conclusion est vraie pour $E(Y_s)$, comme on constate aussitôt; de sorte que, si $s$ est à corps résiduel fini, les valeurs propres du frobenius opérant sur $E^{n-1}(Y_{\bar{s}})$ sont divisibles par \deleted{la} $\operatorname{card}k(t)$ et a fort\deleted{n}iori par $\deleted{\unsure{o|p}}p$. O\unsure{\deleted{r}n|\deleted{n}r} cela est impossible si $D$ estchoisieassez générale, de façon qu'on puisse appliquer Bertini. D'autre part, même si $n$ est impair i.e. $n-1$ pair, de sorte qu'on ne peut plus affirmer que la \deleted{restri}\deleted{b) Bien entendu, ce n'est que par commodité de langage qu'on a} représentation de monodromie sur $E^{n-1}(Y_{\bar t})$ \unsure{\deleted{e}r|r}este simple quand on la restreint à un sous-groupe ouvert quelconque de $G=\operatorname{Gal}(\overline{k(t)}/k(t))$, il sera vrai néanmoins qu'aucun sous-motif non nul de $E^{n-1}(Y_{\bar t})$ n'est de niveau $\handwritten{<}n-1$. En effet, le plus grand sous-motif \inserted{de niveaù $n-1$} de $E^{n-1}(Y_{\bar t})$ \deleted{qui est de niveau $n-1$} définit un sous-espace de $E^{n-1}(Y_{\bar t},\underline{Q}_{\ell})$, qui\deleted{n} est évidemment stable sous le groupe de Galois.

b) Nous avons utilisé le langage des motifs par pure commodité pour l'énoncé de 7.3. Comme ce langage n'est pas justifi\deleted{f}é à l'heure actuelle, il est préférable, pour le moment, de l'énoncer en termes de la filtration par le niveau sur la cohomologie $\ell$-adique, en définissant celle-ci en conformité avec a) (La définition plus naturelle en termes de la filtration de l'espace $Y_{\bar t}$ par les supports devrait être équivalente, mais ce\deleted{v}ci utiliserait la résolution des singularités). Ceci dit, l'énoncé 7.3 (sous la forme plus générale a) ci-dessus) est démontré sans aucune conjecture (pas même Lefschetz vache pour $X$ \ldots).
% archive sheet number: 21
